%% Elevark fra mattelab.no – laget automatisk av bygg.py. Rediger gjerne til eget bruk.
%% Kompiler med lualatex.
\documentclass[11pt,a4paper]{article}
\usepackage[norsk]{babel}
\usepackage[a4paper,margin=2.2cm,top=2cm,bottom=2.4cm]{geometry}
\usepackage{fontspec}
\usepackage{unicode-math}
\setmainfont{TeX Gyre Pagella}
\setsansfont{TeX Gyre Heros}[Scale=MatchLowercase]
\setmathfont{TeX Gyre Pagella Math}
\usepackage{xcolor}
\definecolor{blekk}{HTML}{16213A}
\definecolor{dempet}{HTML}{566077}
\definecolor{rute}{HTML}{CFC9B8}
\definecolor{m1}{HTML}{2F4AB0}
\definecolor{m2}{HTML}{99266D}
\usepackage[most]{tcolorbox}
\usepackage{qrcode}
\usepackage{needspace}
\usepackage{fancyhdr}
\usepackage[hidelinks]{hyperref}
\pdfvariable suppressoptionalinfo \numexpr 1+2+4+8+16+32+64+128+256+512\relax
\setlength{\parindent}{0pt}
\setlength{\parskip}{0.4em}
\pagestyle{fancy}
\fancyhf{}
\renewcommand{\headrulewidth}{0pt}
\fancyfoot[L]{\sffamily\footnotesize\color{dempet}mattelab.no · Polynomdivisjon · elevark uten fasit}
\fancyfoot[R]{\sffamily\footnotesize\color{dempet}Side \thepage}

\newcommand{\fase}[2]{\tcbox[on line, colframe=#2, colback=white, coltext=#2, boxrule=0.6pt, arc=2pt,
  boxsep=0pt, left=3pt, right=3pt, top=1.5pt, bottom=1.5pt, fontupper=\sffamily\bfseries\footnotesize]{#1}}
\newcommand{\steg}[4]{%
  \par\noindent{\color{rute}\rule{\linewidth}{0.4pt}}\par\nopagebreak
  \noindent\begin{minipage}[t]{2.1cm}\vspace{0pt}\fase{#1}{#2}\end{minipage}%
  \begin{minipage}[t]{\dimexpr\linewidth-2.1cm\relax}\vspace{0pt}#3\end{minipage}\par
  \vspace{#4}}
\newcommand{\del}[4]{%
  \par\noindent{\color{rute}\rule{\linewidth}{0.4pt}}\par\nopagebreak
  \noindent\begin{minipage}[t]{1.1cm}\vspace{0pt}\fase{#1}{#2}\end{minipage}%
  \begin{minipage}[t]{\dimexpr\linewidth-1.1cm\relax}\vspace{0pt}#3\end{minipage}\par
  \vspace{#4}}
\newcommand{\oppgave}[2]{\par\Needspace{12\baselineskip}\vspace{0.6em}%
  {\fontsize{14}{17}\selectfont\bfseries{\color{m1}#1}\enspace #2}\par\vspace{0.2em}}

\begin{document}
\begin{minipage}[t]{0.72\linewidth}\vspace{0pt}
{\sffamily\small\color{dempet}Elevark · 1T · mattelab.no}\par\vspace{0.3em}
{\fontsize{24}{28}\selectfont\bfseries Polynomdivisjon}\par\vspace{1em}
{\sffamily\small Navn: \rule{4.6cm}{0.4pt}\quad Klasse: \rule{1.5cm}{0.4pt}\quad Dato: \rule{1.9cm}{0.4pt}}
\end{minipage}\hfill
\begin{minipage}[t]{0.22\linewidth}\vspace{0pt}\raggedleft
\qrcode[height=2.3cm]{https://mattelab.no/utforskere/polynomdivisjon.html}\par\vspace{0.2em}
{\sffamily\scriptsize\color{dempet}Utforskeren}
\end{minipage}

\vspace{1em}
{\sffamily\small\color{dempet}Utforskeren ligger på \url{https://mattelab.no/utforskere/polynomdivisjon.html}}\par
Hver oppgave har tre steg. \textbf{Gjett} først, uten å bruke utforskeren. \textbf{Test} så i utforskeren. \textbf{Forklar} til slutt med egne ord hvorfor det blir slik.

\oppgave{1}{En divisjon som går opp}
\steg{Gjett}{m1}{Regn ut $(x^3 - 2x^2 - 5x + 6) : (x - 1)$ for hånd. Hva blir første ledd i svaret?}{2.2cm}
\steg{Test}{m2}{Sett opp divisjonen. Trykk «Fra start» og gå gjennom stegene. Sammenlign med din egen utregning.}{2.6cm}
\steg{Forklar}{blekk}{Bruk svaret til å faktorisere $x^3 - 2x^2 - 5x + 6$ helt. Hvilke nullpunkter har polynomet?}{3.4cm}

\oppgave{2}{Resten}
\steg{Gjett}{m1}{Hva blir resten når $P(x) = x^2 + 3x + 5$ deles på $x - 2$? Regn ut $P(2)$ også.}{2.2cm}
\steg{Test}{m2}{Sett opp divisjonen. Les av resten og $P(2)$. Endre $r$ og konstantleddet noen ganger. Hva ser du?}{2.6cm}
\steg{Forklar}{blekk}{Start med $P(x) = (x - r) \cdot Q(x) + R$. Forklar hvorfor resten alltid er $P(r)$.}{3.4cm}

\oppgave{3}{Finn et nullpunkt og faktoriser}
\steg{Gjett}{m1}{$P(x) = x^3 - 7x + 6$. Prøv deg fram med $x = \pm 1$, $\pm 2$, $\pm 3$ og $\pm 6$. Hvilke av dem gir $P(x) = 0$?}{2.2cm}
\steg{Test}{m2}{Sett opp $P(x) : (x - 1)$. Legg merke til leddet $0x^2$. Dra $r$ og se når resten blir $0$.}{2.6cm}
\steg{Forklar}{blekk}{Faktoriser $P(x)$ helt. Hvorfor er det nok å prøve divisorene i $6$?}{3.4cm}

\oppgave{4}{En skrå asymptote}
\steg{Gjett}{m1}{Hvordan tror du grafen til $f(x) = \frac{x^2 + x - 1}{x - 1}$ ser ut for store $x$?}{2.2cm}
\steg{Test}{m2}{Sett opp brøken. Les av svaret og resten. Zoom ut og se på grafen og den stiplede linja.}{2.6cm}
\steg{Forklar}{blekk}{Skriv $f(x)$ som et polynom pluss en brøk. Forklar hvorfor grafen nærmer seg linja $y = x + 2$.}{3.4cm}

\oppgave{5}{En ulikhet}
\steg{Gjett}{m1}{For hvilke $x$ er $x^3 - 2x^2 - 5x + 6 > 0$? Bruk faktoriseringen fra oppgave 1.}{2.2cm}
\steg{Test}{m2}{Se på grafen til polynomet i utforskeren for polynomfunksjoner. Slå på fortegnet på x-aksen og sammenlign med svaret ditt.}{2.6cm}
\steg{Forklar}{blekk}{Lag et fortegnsskjema med én linje for hver faktor. Forklar hvordan du finner fortegnet til produktet.}{3.4cm}

\end{document}
